%======================================================================
% GLOSARIO DE ICONOS - SANDRA DESKTOP CONTAINER
% Basado en los iconos FontAwesome reales de la aplicacion
% Usando paquete fontawesome5 de LaTeX
%======================================================================

\chapter*{Glosario de Iconos}
\addcontentsline{toc}{chapter}{Glosario de Iconos}
\markboth{Glosario de Iconos}{Glosario de Iconos}

Este glosario presenta los iconos utilizados en la interfaz de \textbf{Sandra Desktop Container}, organizados por su funcion en la aplicacion.

% Comando para iconos con color sandra, tamaño grande y centrado
\newcommand{\iconocolor}{\color{sandraGreen}\Large}
\newcommand{\iconocent}{\centering\arraybackslash}
\newcommand{\iconocolortext}[1]{{\color{#1}\Large}}

\section*{Iconos de Navegacion Principal}

Iconos del menu lateral (sidebar) de la aplicacion:

\begin{longtable}{>{\centering\arraybackslash}m{1.5cm}p{3.5cm}p{8cm}}
\toprule
\textbf{Icono} & \textbf{Nombre} & \textbf{Funcion} \\
\midrule
\endfirsthead
\toprule
\textbf{Icono} & \textbf{Nombre} & \textbf{Funcion} \\
\midrule
\endhead
\bottomrule
\endfoot

\iconocolor\faIcon{th-large} & \textbf{Dashboard} & Panel principal del sistema con metricas en tiempo real \\
\midrule
\iconocolor\faIcon{network-wired} & \textbf{Conexiones} & Gestion de perfiles de conexion a servidores remotos \\
\midrule
\iconocolor\faIcon{globe} & \textbf{Aplicaciones} & Catalogo y gestion de micro-aplicaciones desplegadas \\
\midrule
\iconocolor\faIcon{shield-alt} & \textbf{Seguridad} & Configuracion de politicas de seguridad y auditoria \\
\midrule
\iconocolor\faIcon{desktop} & \textbf{Monitor} & Visualizacion de metricas de sistema y rendimiento \\
\midrule
\iconocolor\faIcon{shield-alt} & \textbf{Visor Seguro} & Visualizacion protegida de documentos confidenciales \\
\midrule
\iconocolor\faIcon{book} & \textbf{Ayuda Wiki} & Documentacion y guia de usuario integrada \\
\midrule
\iconocolor\faIcon{sliders-h} & \textbf{Configuracion} & Ajustes y preferencias del sistema \\
\midrule
\iconocolor\faIcon{sign-out-alt} & \textbf{Cerrar Sesion} & Cierre seguro de la sesion de usuario \\

\end{longtable}

\section*{Iconos del Dashboard}

Iconos de las tarjetas de metricas en el panel principal:

\begin{longtable}{>{\centering\arraybackslash}m{1.5cm}p{2.5cm}p{9cm}}
\toprule
\textbf{Icono} & \textbf{Metrica} & \textbf{Descripcion} \\
\midrule
\endfirsthead
\toprule
\textbf{Icono} & \textbf{Metrica} & \textbf{Descripcion} \\
\midrule
\endhead
\bottomrule
\endfoot

\iconocolor\faIcon{network-wired} & \textbf{Red} & Estado de conexion local, direccion IP, latencia y throughput \\
\midrule
\iconocolor\faIcon{hdd} & \textbf{Espacio} & Capacidad disponible en disco para contenedores y datos \\
\midrule
\iconocolor\faIcon{database} & \textbf{Core DB} & Tamano y estado de la base de datos SQLite interna cifrada \\

\end{longtable}

\section*{Iconos de Estado}

Iconos que indican el estado de elementos o alertas:

\begin{longtable}{>{\centering\arraybackslash}m{1.5cm}p{3.5cm}p{7.5cm}}
\toprule
\textbf{Icono} & \textbf{Estado} & \textbf{Significado} \\
\midrule
\endfirsthead
\toprule
\textbf{Icono} & \textbf{Estado} & \textbf{Significado} \\
\midrule
\endhead
\bottomrule
\endfoot

{\color{green}\faIcon{check-circle}} & \textbf{Exito} & Operacion completada correctamente, conexion estable \\
\midrule
{\color{red}\faIcon{times-circle}} & \textbf{Error} & Fallo en operacion, error critico o conexion rechazada \\
\midrule
{\color{orange}\faIcon{exclamation-triangle}} & \textbf{Advertencia} & Alerta que requiere atencion del usuario \\
\midrule
{\color{blue}\faIcon{info-circle}} & \textbf{Informacion} & Detalles adicionales o ayuda contextual \\
\midrule
\iconocolor\faIcon{globe-americas} & \textbf{Modo Libre} & Aplicacion en modo de navegacion directa sin proxy \\
\midrule
{\color{green}\faIcon{shield-alt}} & \textbf{Proxy Activo} & Conexion protegida mediante proxy seguro \\

\end{longtable}

\section*{Iconos de Accion}

Iconos para botones de accion en la interfaz:

\begin{longtable}{>{\centering\arraybackslash}m{1.5cm}p{4cm}p{7cm}}
\toprule
\textbf{Icono} & \textbf{Accion} & \textbf{Funcion} \\
\midrule
\endfirsthead
\toprule
\textbf{Icono} & \textbf{Accion} & \textbf{Funcion} \\
\midrule
\endhead
\bottomrule
\endfoot

\iconocolor\faIcon{plus} & \textbf{Nuevo/Agregar} & Crear nuevo elemento (conexion, aplicacion, etc.) \\
\midrule
\iconocolor\faIcon{minus} & \textbf{Quitar} & Eliminar elemento seleccionado de una lista \\
\midrule
\iconocolor\faIcon{edit} & \textbf{Editar} & Modificar configuracion o propiedades \\
\midrule
\iconocolor\faIcon{save} & \textbf{Guardar} & Persistir cambios en configuracion o datos \\
\midrule
\iconocolor\faIcon{trash} & \textbf{Eliminar} & Borrado seguro con confirmacion \\
\midrule
\iconocolor\faIcon{folder-minus} & \textbf{Desactivar} & Quitar aplicacion sin eliminar datos \\
\midrule
\iconocolor\faIcon{sync} & \textbf{Sincronizar} & Actualizar estado con servidor remoto \\
\midrule
\iconocolor\faIcon{sync-alt} & \textbf{Actualizar} & Buscar actualizaciones de aplicaciones \\
\midrule
\iconocolor\faIcon{redo-alt} & \textbf{Refrescar} & Recargar contenido actual \\
\midrule
\iconocolor\faIcon{search} & \textbf{Buscar} & Funcion de busqueda en documentos o registros \\
\midrule
\iconocolor\faIcon{filter} & \textbf{Filtrar} & Aplicar filtros a listas o resultados \\
\midrule
\iconocolor\faIcon{download} & \textbf{Descargar} & Recibir datos o archivos del servidor \\
\midrule
\iconocolor\faIcon{upload} & \textbf{Subir} & Enviar datos o archivos al servidor \\
\midrule
\iconocolor\faIcon{arrow-right} & \textbf{Abrir/Siguiente} & Abrir aplicacion o navegar adelante \\
\midrule
\iconocolor\faIcon{arrow-left} & \textbf{Anterior} & Navegar hacia atras o cancelar \\
\midrule
\iconocolor\faIcon{expand} & \textbf{Expandir} & Ampliar vista o mostrar detalles \\
\midrule
\iconocolor\faIcon{compress} & \textbf{Comprimir} & Reducir vista o colapsar panel \\
\midrule
\iconocolor\faIcon{power-off} & \textbf{Apagar} & Cierre seguro del sistema \\
\midrule
\iconocolor\faIcon{times} & \textbf{Cerrar} & Cerrar pestana o dialogo \\

\end{longtable}

\section*{Iconos de Seguridad}

Iconos relacionados con seguridad y control de acceso:

\begin{longtable}{>{\centering\arraybackslash}m{1.5cm}p{4cm}p{7cm}}
\toprule
\textbf{Icono} & \textbf{Elemento} & \textbf{Descripcion} \\
\midrule
\endfirsthead
\toprule
\textbf{Icono} & \textbf{Elemento} & \textbf{Descripcion} \\
\midrule
\endhead
\bottomrule
\endfoot

\iconocolor\faIcon{lock} & \textbf{Bloqueado} & Recurso protegido o documento encriptado \\
\midrule
\iconocolor\faIcon{unlock} & \textbf{Desbloqueado} & Acceso autorizado o documento desencriptado \\
\midrule
\iconocolor\faIcon{key} & \textbf{Cifrado} & Gestion de claves y configuracion criptografica \\
\midrule
\iconocolor\faIcon{fingerprint} & \textbf{Identidad} & Client ID unico del dispositivo vinculado \\
\midrule
\iconocolor\faIcon{user} & \textbf{Usuario} & Cuenta del operador actual \\
\midrule
\iconocolor\faIcon{users} & \textbf{Area/Departamento} & Unidad organizativa asignada \\
\midrule
\iconocolor\faIcon{link} & \textbf{Proxy/Ruta} & Punto de acceso seguro configurado \\
\midrule
\iconocolor\faIcon{server} & \textbf{Servidor} & Nodo servidor en el ecosistema SDC \\
\midrule
\iconocolor\faIcon{plug} & \textbf{Conexion} & Perfil de conexion activa o inactiva \\

\end{longtable}

\section*{Iconos de Hardware}

Iconos de componentes hardware y recursos del sistema:

\begin{longtable}{>{\centering\arraybackslash}m{1.5cm}p{3cm}p{8cm}}
\toprule
\textbf{Icono} & \textbf{Componente} & \textbf{Indicador} \\
\midrule
\endfirsthead
\toprule
\textbf{Icono} & \textbf{Componente} & \textbf{Indicador} \\
\midrule
\endhead
\bottomrule
\endfoot

\iconocolor\faIcon{microchip} & \textbf{CPU/Procesador} & Uso porcentual y nucleos activos \\
\midrule
\iconocolor\faIcon{memory} & \textbf{RAM/Memoria} & Memoria disponible y uso actual \\
\midrule
\iconocolor\faIcon{wifi} & \textbf{WiFi} & Conexion inalambrica activa \\
\midrule
\iconocolor\faIcon{ethernet} & \textbf{Ethernet} & Conexion cableada de red \\

\end{longtable}

\section*{Iconos de Contenido}

Iconos para tipos de aplicaciones y asistentes:

\begin{longtable}{>{\centering\arraybackslash}m{1.5cm}p{4cm}p{7cm}}
\toprule
\textbf{Icono} & \textbf{Tipo} & \textbf{Descripcion} \\
\midrule
\endfirsthead
\toprule
\textbf{Icono} & \textbf{Tipo} & \textbf{Descripcion} \\
\midrule
\endhead
\bottomrule
\endfoot

\iconocolor\faIcon{code} & \textbf{Aplicacion} & Micro-aplicacion o servicio web desplegado \\
\midrule
\iconocolor\faIcon{file-alt} & \textbf{Documento} & Archivo gestionado por el sistema \\
\midrule
\iconocolor\faIcon{eye} & \textbf{Visor} & Modo de visualizacion segura de documentos \\
\midrule
\iconocolor\faIcon{robot} & \textbf{Sandra AI} & Asistente de inteligencia artificial integrado \\
\midrule
\iconocolor\faIcon{history} & \textbf{Historial} & Registro de auditoria y logs de actividad \\

\end{longtable}

\section*{Colores de Estado}

La interfaz utiliza el siguiente codigo de colores para indicar estados:

\begin{table}[H]
\centering
\begin{tabular}{>{\centering\arraybackslash}p{3cm}p{3cm}p{6cm}}
\toprule
\textbf{Color} & \textbf{Estado} & \textbf{Significado} \\
\midrule
\cellcolor{sandraGreen} & \textbf{Verde} & Operacion exitosa, conexion activa, sistema estable \\
\midrule
\cellcolor{sandraAccent} & \textbf{Ambar} & Advertencia, atencion requerida, proceso en curso \\
\midrule
\cellcolor{red!70} & \textbf{Rojo} & Error critico, conexion fallida, alerta de seguridad \\
\midrule
\cellcolor{blue!50} & \textbf{Azul} & Informacion, estado neutro, proceso activo \\
\midrule
\cellcolor{gray!50} & \textbf{Gris} & Deshabilitado, inactivo, no disponible \\
\bottomrule
\end{tabular}
\caption{Codificacion de colores de estado en SDC}
\label{tab:colores-estado}
\end{table}

\section*{Abreviaturas y Acronimos}

\begin{longtable}{p{3cm}p{10cm}}
\toprule
\textbf{Acronimo} & \textbf{Significado} \\
\midrule
\endfirsthead
\toprule
\textbf{Acronimo} & \textbf{Significado} \\
\midrule
\endhead
\bottomrule
\endfoot

SDC & Sandra Desktop Container \\
SID & Secure ID (Identificador Seguro) \\
IPC & Inter-Process Communication (Comunicacion entre procesos) \\
WSS & WebSocket Secure (WebSocket con TLS) \\
TLS & Transport Layer Security (Seguridad de capa de transporte) \\
AES & Advanced Encryption Standard (Estandar de cifrado avanzado) \\
GCM & Galois/Counter Mode (Modo de operacion de cifrado) \\
GDPR & General Data Protection Regulation (Reglamento general de proteccion de datos) \\
LOPD & Ley Organica de Proteccion de Datos \\
CPU & Central Processing Unit (Unidad central de procesamiento) \\
RAM & Random Access Memory (Memoria de acceso aleatorio) \\
OS & Operating System (Sistema operativo) \\
VPN & Virtual Private Network (Red privada virtual) \\
PoLP & Principle of Least Privilege (Principio de minimo privilegio) \\
ISO & International Organization for Standardization \\
IEC & International Electrotechnical Commission \\
NIST & National Institute of Standards and Technology \\

\end{longtable}
